\section{格子ボルツマン法の基礎理論}
\label{sec:lbm}

\subsection{ボルツマン方程式とBGK近似}

格子ボルツマン法（Lattice Boltzmann Method, LBM）は，Navier-Stokes方程式を直接離散化するのではなく，より基礎的な運動論レベルの方程式——ボルツマン輸送方程式——から出発し，Chapman-Enskog展開によってマクロスケールで圧縮性の弱いNavier-Stokes方程式を導出するアプローチである。

分子速度分布関数 \(f(\vec x,\vec\xi,t)\) に対するボルツマン方程式を，単一緩和時間（BGK）近似で時間・空間・速度空間すべてについて離散化すると次の格子ボルツマン方程式を得る：
\begin{align}
	f_i(\vec x+\vec e_i\Delta t,\ t+\Delta t) - f_i(\vec x,t)
	= -\frac{1}{\tau}\left[f_i(\vec x,t) - f_i^{eq}(\vec x,t)\right],
	\label{eq:lbe}
\end{align}
ここで \(f_i\) は離散速度方向 \(i\) の分布関数，\(\vec e_i\) は格子速度ベクトル，\(\tau\) は無次元緩和時間，\(f_i^{eq}\) はMaxwell-Boltzmann分布の低マッハ数展開である局所平衡分布である。緩和時間は動粘性係数と
\begin{align}
	\nu = c_s^2\left(\tau - \frac{1}{2}\right)\Delta t
	\label{eq:viscosity}
\end{align}
の関係で結びつく。ここで \(c_s = 1/\sqrt3\)（格子単位）は格子音速である。\eref{lbe}の右辺で表される衝突ステップと，左辺の並進（ストリーミング）ステップはいずれも局所演算のみで構成され，格子点間の通信は最近傍への値のコピーに限られる。この局所性が，GPUによる大規模並列化との親和性の根拠である。

\subsection{離散速度セット}

本プロジェクトでは2つの離散速度セットを用いる。

\begin{itemize}
	\item \textbf{D3Q19}：標準的な3次元計算に用いられるモデルで，メモリ効率と精度のバランスが良い。デバッグおよび低レイノルズ数での検証に用いる。
	\item \textbf{D3Q27}：本番計算（縫い目球・乱流）に用いる既定モデル。\secref{cumulant}で述べる中心モーメント衝突演算子は，D3Q27を必須とする——3次モーメント \(\sum_q w_q c_{x}^2c_{y}^2c_{z}^2\) がD3Q19では恒等的に0となる一方，D3Q27では \(c_s^6\) となり，3方向すべてに独立に分解できるのはD3Q27のみだからである。
\end{itemize}

本プロジェクトが対象とするレイノルズ数（\(Re\sim2\times10^5\)）は乱流域であり数値安定性が課題となるため，本番計算はD3Q27と中心モーメント衝突演算子の組を推奨構成とし，D3Q19+BGKは検証・軽量計算用の代替構成として残す。単一GPU（8\,GB VRAM）という制約下では，D3Q27はD3Q19よりメモリ使用量が増える点にも留意し，解像度とのトレードオフを\secref{resolution-budget}で検討する。

\subsection{格子単位系と物理量の変換}
\label{sec:lattice-units}

LBMは無次元格子単位（\(\Delta x=\Delta t=1\)）で計算するため，物理パラメータ（音速比・粘性・球速）を格子単位に変換する層が必要になる。設計上の要件は次の2点である。

\begin{itemize}
	\item 格子マッハ数 \(Ma_{lattice}=U_{lattice}/c_{s,lattice}\) を0.05--0.1程度に保つよう，代表速度と格子解像度から \(\Delta t\) を自動決定する。密結合方式（\secref{coupling}）では球速が飛行中にわずかに変化するため，初期速度を基準にマージンを持たせて設定する。
	\item 空間解像度 \(\Delta x\) はボール直径を何格子点で解像するか（\(N=D/\Delta x\)）で規定する。既定値は縫い目高さと境界層厚を何格子点で捉えるかという要件から決定する（\secref{resolution-budget}）。
\end{itemize}

一様な格子マッハ数で解像度を切る限り，格子粘性 \(\nu_{lattice}=\nu\,\Delta t/\Delta x^2\) は解像度 \(N\) に比例して大きくなる。したがって\emph{格子を細かくするほど \(\tau\) は\(1/2\)から離れて安定側に動く}——本プロジェクトの代表条件では，40点/Dで \(\tau-1/2=3.1\times10^{-5}\)，80点/Dで \(6.2\times10^{-5}\) となる。すなわち本番計算のコストを決めているのは数値安定性ではなく，飛行時間全体に必要なタイムステップ数である。
